\section{Teoremas e identidades para reconocer soluciones}
\whenuse{Antes de programar, cuando el problema pide contar, demostrar una cota,
trabajar m\'odulo primo o transformar una suma aparentemente cuadr\'atica.}
\complexity{Las f\'ormulas son $O(1)$ una vez conocidos sus datos;
inclusi\'on--exclusi\'on cuesta normalmente $O(2^k)$ y Burnside recorre las
simetr\'ias. Precomputaciones y factorizaciones conservan sus propios costos.}

\subsection{Conteo y cotas}
\begin{itemize}
  \item \textbf{Palomar:} al repartir $m$ objetos entre $n$ cajas, alguna recibe
  al menos $\lceil m/n\rceil$. Se usa para demostrar que una repetici\'on o
  colisi\'on es inevitable sin identificar cu\'al. Por ejemplo, entre $13$
  personas al menos dos cumplen a\~nos en el mismo mes.
  \item \textbf{Handshake:}
  \[
      \sum_v deg(v)=2|E|.
  \]
  Cada arista aporta uno al grado de cada extremo, por eso se cuenta dos veces.
  Como la suma es par, la cantidad de v\'ertices de grado impar tambi\'en es par.
  Es \'util para validar secuencias de grados y recorridos eulerianos.
  \item \textbf{Suma triangular:}
  \[
      1+2+\cdots+k=\frac{k(k+1)}2.
  \]
  Aparece cuando la iteraci\'on $i$ cuesta $i$, se eligen pares o cada paso
  consume una unidad m\'as. Si la suma no supera $n$, entonces
  $k\approx\sqrt{2n}$ y por tanto $k=O(\sqrt n)$.
  \item \textbf{Serie arm\'onica entera:}
  \[
      \sum_{i=1}^n\left\lfloor\frac ni\right\rfloor=O(n\log n).
  \]
  Cuenta, por ejemplo, todos los m\'ultiplos procesados por bucles de la forma
  \texttt{for(j=i;j<=n;j+=i)}. Adem\'as, $\lfloor n/i\rfloor$ toma s\'olo
  $O(\sqrt n)$ valores distintos; para un cociente $q$, el \'ultimo \'indice con
  el mismo valor es $\lfloor n/q\rfloor$, permitiendo agrupar rangos.
  \item \textbf{Inclusi\'on--exclusi\'on:} suma los casos que cumplen cada
  propiedad y corrige los contados varias veces:
  \[
      |A\cup B|=|A|+|B|-|A\cap B|,
  \]
  \[
      |A\cup B\cup C|=\sum|A_i|-\sum|A_i\cap A_j|+|A\cap B\cap C|.
  \]
  Para $k$ propiedades se recorren normalmente las $2^k$ intersecciones. Se usa
  para contar objetos que cumplen al menos una restricci\'on o, tomando el
  complemento, los que no cumplen ninguna.
\end{itemize}

\subsection{Combinatoria}
\begin{itemize}
  \item \textbf{Reglas de suma y producto:} si los casos son excluyentes, sus
  cantidades se suman. Si una elecci\'on de $a$ formas va seguida de otra
  independiente de $b$ formas, se multiplican: $ab$. Por ejemplo, $3$ camisas y
  $2$ pantalones producen $3\cdot2=6$ conjuntos.
  \item \textbf{Permutaciones:} ordenar $n$ objetos distintos produce $n!$
  \'ordenes. Elegir y ordenar $k$ de ellos produce
  \[
      P(n,k)=n(n-1)\cdots(n-k+1)=\frac{n!}{(n-k)!}.
  \]
  Se usa cuando el orden importa; elegir presidente y vicepresidente entre $n$
  personas da $P(n,2)$, no $\binom n2$.
  \item \textbf{Combinaciones:} elegir $k$ objetos entre $n$ sin importar el
  orden:
  \[
      \binom nk=\frac{n!}{k!(n-k)!}=\binom n{n-k}.
  \]
  Si un multiconjunto tiene cantidades $c_1,\ldots,c_r$ con $\sum c_i=n$, sus
  \'ordenes distintos son
  \[
      \frac{n!}{c_1!c_2!\cdots c_r!}.
  \]
  Por ejemplo, \texttt{AAB} tiene $3!/2!=3$ permutaciones distintas.
  \item \textbf{Subconjuntos:} cada uno de los $n$ elementos puede estar o no,
  por lo que existen $2^n$ subconjuntos. Los de tama\~no $k$ son $\binom nk$ y
  \[
      \sum_{k=0}^n\binom nk=2^n.
  \]
  \item \textbf{Identidad de Pascal:}
  \[
      \binom nk=\binom{n-1}{k-1}+\binom{n-1}{k}.
  \]
  Al fijar un elemento, una selecci\'on o lo contiene o no lo contiene. Permite
  calcular combinaciones mediante DP, incluso cuando no existe divisi\'on
  modular.
  \item \textbf{Hockey-stick:}
  \[
      \sum_{i=r}^n\binom ir=\binom{n+1}{r+1}.
  \]
  Convierte una suma de combinaciones consecutivas en una sola. Combinatoriamente,
  clasifica cada selecci\'on de $r+1$ elementos seg\'un cu\'al es su elemento
  mayor.
  \item \textbf{Binomio de Newton:}
  \[
      (x+y)^n=\sum_{k=0}^n\binom nkx^ky^{n-k}.
  \]
  El coeficiente $\binom nk$ cuenta en cu\'ales $k$ de los $n$ factores se elige
  $x$. Sustituir valores produce identidades; con $x=y=1$ se obtiene
  $\sum_k\binom nk=2^n$.
  \item \textbf{Stars and bars:} la cantidad de soluciones enteras no negativas
  de
  \[
      x_1+\cdots+x_k=n
  \]
  es $\binom{n+k-1}{k-1}$: se ordenan $n$ estrellas y $k-1$ separadores. Si
  todas las variables deben ser positivas, entregar una unidad a cada una deja
  $n-k$ por repartir y resulta $\binom{n-1}{k-1}$. Sirve para distribuir objetos
  id\'enticos entre cajas distinguibles.
  \item \textbf{Vandermonde:}
  \[
      \sum_i\binom ri\binom s{n-i}=\binom{r+s}n.
  \]
  Elegir $n$ elementos de dos grupos de tama\~nos $r$ y $s$ puede hacerse
  directamente o fijando cu\'antos $i$ se toman del primero.
  \item \textbf{Catalan:}
  \[
      C_n=\frac1{n+1}\binom{2n}n
         =\binom{2n}n-\binom{2n}{n+1}.
  \]
  Cuenta estructuras que nunca cruzan una frontera: par\'entesis balanceados con
  $n$ pares, caminos que no bajan de cero, triangulaciones de un pol\'igono de
  $n+2$ lados y varios \'arboles ordenados. Los primeros valores son
  $1,1,2,5,14,42$.
  \item \textbf{Desarreglos:} permutaciones donde ning\'un elemento queda en su
  posici\'on original:
  \[
      D_n=(n-1)(D_{n-1}+D_{n-2}),\qquad D_0=1,\ D_1=0.
  \]
  Al decidir d\'onde va el primer elemento, el caso restante se reduce a un
  desarreglo de $n-1$ o de $n-2$. Aparece en asignaciones donde nadie puede
  recibir su objeto original.
  \item \textbf{Burnside:} para contar configuraciones iguales bajo un grupo de
  rotaciones o reflexiones,
  \[
      \#\text{clases}=\frac1{|G|}\sum_{g\in G}|\operatorname{Fix}(g)|.
  \]
  Para cada simetr\'ia se cuentan las configuraciones que no cambian y luego se
  promedia. Es \'util para collares, pulseras, coloraciones y tableros considerados
  iguales bajo simetr\'ias.
\end{itemize}

\subsection{Teor\'ia de n\'umeros}
Si $n=\prod p_i^{a_i}$, cada divisor elige para $p_i$ un exponente entre $0$ y
$a_i$. De esta observaci\'on se obtienen:
\[
    \tau(n)=\prod_i(a_i+1),\qquad
    \sigma(n)=\prod_i\frac{p_i^{a_i+1}-1}{p_i-1},\qquad
    \varphi(n)=n\prod_{p\mid n}\left(1-\frac1p\right).
\]
$\tau(n)$ cuenta divisores, $\sigma(n)$ suma sus valores y $\varphi(n)$ cuenta
enteros de $1$ a $n$ coprimos con $n$. Por ejemplo,
$12=2^2\cdot3$ da $\tau(12)=6$, $\sigma(12)=28$ y $\varphi(12)=4$.
\begin{itemize}
  \item \textbf{B\'ezout:} existen enteros $x,y$ tales que
  \[
      ax+by=\gcd(a,b).
  \]
  Por tanto, $ax+by=c$ tiene soluciones enteras si y s\'olo si
  $\gcd(a,b)\mid c$. Euclides extendido encuentra $x,y$ y permite construir
  inversos modulares cuando el gcd es $1$.
  \item \textbf{Peque\~no teorema de Fermat:} si $p$ es primo y $p\nmid a$,
  \[
      a^{p-1}\equiv1\pmod p.
  \]
  Entonces $a^{-1}\equiv a^{p-2}\pmod p$ y los exponentes pueden reducirse
  m\'odulo $p-1$ cuando la base no es divisible por $p$.
  \item \textbf{Teorema de Euler:} si $\gcd(a,m)=1$, entonces
  \[
      a^{\varphi(m)}\equiv1\pmod m.
  \]
  Generaliza Fermat a m\'odulos no primos y permite reducir exponentes m\'odulo
  $\varphi(m)$ bajo la condici\'on de coprimalidad.
  \item \textbf{Wilson:} $p>1$ es primo si y s\'olo si
  \[
      (p-1)!\equiv-1\pmod p.
  \]
  Es \'util en demostraciones y algunos c\'alculos cercanos a $p$, pero calcular
  el factorial cuesta $O(p)$ y no es una prueba de primalidad eficiente para
  valores grandes.
\end{itemize}

\subsection{Grafos y geometr\'ia}
\begin{itemize}
  \item \textbf{Caracterizaci\'on de \'arboles:} un \'arbol conectado con $n$
  v\'ertices tiene exactamente $n-1$ aristas. En un \'arbol existe un \'unico
  camino entre cada par; agregar una arista crea exactamente un ciclo y quitar
  cualquiera lo desconecta. Un grafo conectado con $n-1$ aristas es un \'arbol.
  \item \textbf{Caminos eulerianos:} ignorando v\'ertices aislados, las aristas
  deben pertenecer a una sola componente conexa. Existe circuito euleriano si
  todos los grados son pares; existe camino abierto que usa cada arista una vez
  si exactamente dos grados son impares, que ser\'an sus extremos. Con cero
  impares tambi\'en existe camino, pero comienza y termina en el mismo v\'ertice.
  \item \textbf{F\'ormula de Euler planar:} para un grafo planar conexo,
  \[
      V-E+F=2,
  \]
  donde $F$ incluye la cara exterior. Para un grafo planar simple con $V\ge3$
  implica $E\le3V-6$, una cota \'util para demostrar que un grafo no es planar o
  que un algoritmo sobre todas sus aristas es lineal en $V$.
  \item \textbf{Teorema de Pick:} para un pol\'igono simple cuyos v\'ertices son
  puntos enteros,
  \[
      A=I+\frac B2-1,
  \]
  donde $I$ es la cantidad de puntos enteros interiores y $B$ los de la frontera.
  Los puntos de frontera de un lado $(x_1,y_1)$--$(x_2,y_2)$ se cuentan mediante
  $\gcd(|x_2-x_1|,|y_2-y_1|)$ al sumar todos los lados.
\end{itemize}
